\documentclass{article}
\usepackage[utf8]{inputenc}
\usepackage{amsmath,amssymb}
\usepackage[most]{tcolorbox}
\usepackage{geometry}
\geometry{margin=2.5cm}

\newcommand{\step}[1]{\par\noindent\hangindent=3.5em\hangafter=1%
  \makebox[3.5em][l]{\textbf{#1.}}}
\newcommand{\steptag}[1]{\hfill{\small\textsf{[#1]}}}

\newtcolorbox{proofbox}[1]{
  colback=gray!5, colframe=gray!60,
  fonttitle=\bfseries, title={#1},
  breakable, enhanced,
  left=4pt, right=4pt, top=4pt, bottom=4pt,
  before skip=10pt, after skip=10pt
}

\begin{document}

\begin{proofbox}{Residual distance bound: let C be the convex hull of fini...}
\step{1} Residual distance bound: let C be the convex hull of finitely many points of R\textasciicircum{}n, let b\_1..b\_M, c\_1..c\_N >= 0 with m = sum\_j b\_j - sum\_k c\_k > 0, let p\_j, r\_k be points of R\textasciicircum{}n with q = (sum\_j b\_j p\_j - sum\_k c\_k r\_k) / m, and let D\_k >= 0 satisfy ||x - r\_k||\_1 <= D\_k for all x in C and each k; then m * dist\_1(q, C) <= sum\_j b\_j * dist\_1(p\_j, C) + sum\_k c\_k * D\_k. \steptag{validated} \\[6pt]
\step{1.1} Because C is a nonempty finite convex hull in R\textasciicircum{}n, it is compact; hence for every j there exists beta\_j in C with ||p\_j - beta\_j||\_1 = dist\_1(p\_j, C). \steptag{validated} \\[6pt]
\step{1.1.1} A nonempty convex hull of finitely many points in R\textasciicircum{}n is compact, and for each fixed p\_j the map y -> ||p\_j - y||\_1 is continuous on C. \steptag{validated} \\[4pt]
\step{1.1.2} A continuous real-valued function on compact C attains its minimum, so choose beta\_j in C minimizing y -> ||p\_j - y||\_1; by the definition of dist\_1(p\_j, C) as this infimum, ||p\_j - beta\_j||\_1 = dist\_1(p\_j, C). \steptag{validated} \\[4pt]
\step{1.2} Let B = sum\_j b\_j and G = sum\_k c\_k. Since m = B - G > 0 and all c\_k >= 0, B = m + G > 0; therefore beta = (sum\_j b\_j beta\_j) / B is a convex combination of points of C and so beta lies in C. \steptag{validated} \\[6pt]
\step{1.2.1} Because every c\_k is nonnegative, G = sum\_k c\_k is nonnegative; from m = B - G and m > 0 we get B = m + G > 0. \steptag{validated} \\[4pt]
\step{1.2.2} For each j, lambda\_j = b\_j / B is nonnegative because b\_j >= 0 and B > 0, and sum\_j lambda\_j = (sum\_j b\_j) / B = 1. \steptag{validated} \\[4pt]
\step{1.2.3} The point beta = sum\_j lambda\_j beta\_j is therefore a convex combination of the points beta\_j in C; since C is convex, beta lies in C. \steptag{validated} \\[4pt]
\step{1.3} With beta as above and q = (sum\_j b\_j p\_j - sum\_k c\_k r\_k) / m, the vector identity m*(q - beta) = sum\_j b\_j*(p\_j - beta\_j) + sum\_k c\_k*(beta - r\_k) holds. \steptag{validated} \\[6pt]
\step{1.3.1} Multiplying the definition of q by m gives m*q = sum\_j b\_j*p\_j - sum\_k c\_k*r\_k, so m*(q - beta) = sum\_j b\_j*p\_j - sum\_k c\_k*r\_k - m*beta. \steptag{validated} \\[4pt]
\step{1.3.2} From B = sum\_j b\_j and beta = (sum\_j b\_j*beta\_j)/B we have B*beta = sum\_j b\_j*beta\_j; from G = sum\_k c\_k we have G*beta = sum\_k c\_k*beta. \steptag{validated} \\[4pt]
\step{1.3.3} Since m = B - G, m*beta = B*beta - G*beta = sum\_j b\_j*beta\_j - sum\_k c\_k*beta. \steptag{validated} \\[6pt]
\step{1.3.3.1} In the current setup, B = sum\_j b\_j, G = sum\_k c\_k, and the root hypothesis gives m = sum\_j b\_j - sum\_k c\_k; hence m = B - G, so scalar distributivity gives m*beta = (B - G)*beta = B*beta - G*beta. \steptag{validated} \\[4pt]
\step{1.3.3.2} The validated node 1.3.2 supplies the identities B*beta = sum\_j b\_j*beta\_j and G*beta = sum\_k c\_k*beta. \steptag{validated} \\[4pt]
\step{1.3.3.3} Combining 1.3.3.1 with the two identities from 1.3.3.2, substitute B*beta = sum\_j b\_j*beta\_j and G*beta = sum\_k c\_k*beta into m*beta = B*beta - G*beta to obtain m*beta = sum\_j b\_j*beta\_j - sum\_k c\_k*beta, which is the asserted equality chain. \steptag{validated} \\[6pt]
\step{1.3.3.3.1} The dependency-scoping premise is now satisfied: the declared dependencies 1.3.3.1 and 1.3.3.2 are validated in this workspace. Node 1.3.3.1 supplies m*beta = B*beta - G*beta, and node 1.3.3.2 supplies B*beta = sum\_j b\_j*beta\_j and G*beta = sum\_k c\_k*beta. Substituting these two validated identities into the validated equality m*beta = B*beta - G*beta gives m*beta = sum\_j b\_j*beta\_j - sum\_k c\_k*beta, with no appeal to an unvalidated premise. \steptag{archived} \\[4pt]
\step{1.3.3.3.2} Because dependencies 1.3.3.1 and 1.3.3.2 are validated and are declared as the only inputs for this node, the substitution is now dependency-scoped: 1.3.3.1 gives m*beta = B*beta - G*beta, while 1.3.3.2 gives B*beta = sum\_j b\_j*beta\_j and G*beta = sum\_k c\_k*beta; replacing B*beta and G*beta in the first equality by these equal vectors yields m*beta = sum\_j b\_j*beta\_j - sum\_k c\_k*beta. \steptag{validated} \\[4pt]
\step{1.3.4} Substituting this expression for m*beta into m*(q - beta) and regrouping gives sum\_j b\_j*(p\_j - beta\_j) + sum\_k c\_k*(beta - r\_k). \steptag{validated} \\[4pt]
\step{1.4} The positive-part residual satisfies ||sum\_j b\_j*(p\_j - beta\_j)||\_1 <= sum\_j b\_j * dist\_1(p\_j, C). \steptag{validated} \\[6pt]
\step{1.4.1} By the triangle inequality for the l1 norm, ||sum\_j b\_j*(p\_j - beta\_j)||\_1 <= sum\_j ||b\_j*(p\_j - beta\_j)||\_1. \steptag{validated} \\[4pt]
\step{1.4.2} Because b\_j >= 0, absolute homogeneity gives ||b\_j*(p\_j - beta\_j)||\_1 = b\_j*||p\_j - beta\_j||\_1 for every j. \steptag{validated} \\[4pt]
\step{1.4.3} By the chosen beta\_j, ||p\_j - beta\_j||\_1 = dist\_1(p\_j, C) for every j; substituting this into the preceding sum gives the stated bound. \steptag{validated} \\[4pt]
\step{1.5} The negative-part residual satisfies ||sum\_k c\_k*(beta - r\_k)||\_1 <= sum\_k c\_k * D\_k. \steptag{validated} \\[6pt]
\step{1.5.1} By the triangle inequality and absolute homogeneity for the l1 norm, using c\_k >= 0, ||sum\_k c\_k*(beta - r\_k)||\_1 <= sum\_k c\_k*||beta - r\_k||\_1. \steptag{validated} \\[4pt]
\step{1.5.2} Since beta lies in C and the hypotheses give ||x - r\_k||\_1 <= D\_k for every x in C and each k, taking x = beta yields ||beta - r\_k||\_1 <= D\_k for every k. \steptag{validated} \\[6pt]
\step{1.5.2.1} Within this proof, beta in C is supplied by the earlier beta-in-C step 1.2: B = sum\_j b\_j is positive, beta = (sum\_j b\_j beta\_j)/B = sum\_j (b\_j/B) beta\_j, the coefficients b\_j/B are nonnegative and sum to 1, each beta\_j lies in C, and C is convex; hence beta lies in C. \steptag{validated} \\[6pt]
\step{1.5.2.1.1} By validated node 1.2, beta lies in C. That node proves the needed convex-combination fact from B=sum\_j b\_j, G=sum\_k c\_k, m=B-G>0, c\_k>=0, b\_j>=0, the chosen beta\_j in C, and convexity of C; hence the present node may use beta in C through its declared dependency on 1.2. \steptag{archived} \\[4pt]
\step{1.5.2.1.2} By validated node 1.2, beta lies in C. Node 1.2 proves the needed convex-combination fact: B=sum\_j b\_j and G=sum\_k c\_k satisfy B=m+G>0, the coefficients b\_j/B are nonnegative and sum to 1, and beta=sum\_j (b\_j/B) beta\_j is a convex combination of points beta\_j in C. Since C is convex, beta lies in C, so the beta-in-C premise used by 1.5.2.1 is supplied by this declared validated dependency. \steptag{validated} \\[4pt]
\step{1.5.2.2} For each fixed k, the root hypothesis on D\_k is the universal statement: for every x in C, ||x - r\_k||\_1 <= D\_k. Applying that statement to the particular admissible point x = beta, using beta in C from 1.5.2.1, gives ||beta - r\_k||\_1 <= D\_k. Since k was arbitrary, the inequality holds for every k. \steptag{validated} \\[6pt]
\step{1.5.2.2.1} The admissibility premise beta in C is supplied here by the validated dependency 1.2: node 1.2 proves that, with B = sum\_j b\_j > 0 and beta = (sum\_j b\_j beta\_j)/B, beta is a convex combination of points beta\_j in C, hence beta lies in C. This child records that beta in C is in scope for the universal-instantiation step below. \steptag{validated} \\[6pt]
\step{1.5.2.2.1.1} Dependency-scoped repair for ch-390c4b1b8facde9b: this child depends directly on validated node 1.2. Node 1.2 proves that B=sum\_j b\_j and G=sum\_k c\_k satisfy B=m+G>0, that beta=(sum\_j b\_j beta\_j)/B is the convex combination sum\_j (b\_j/B) beta\_j of points beta\_j in C, and hence beta lies in C. Therefore the beta-in-C premise asserted by node 1.5.2.2.1 is now supplied through an actual logical dependency, not merely through validation gating. \steptag{validated} \\[4pt]
\step{1.5.2.2.2} Fix an index k. The root hypothesis on D\_k says that every x in C satisfies ||x - r\_k||\_1 <= D\_k. By 1.5.2.2.1, beta lies in C; substituting x = beta into that universal statement therefore gives ||beta - r\_k||\_1 <= D\_k for this fixed k. \steptag{validated} \\[6pt]
\step{1.5.2.2.2.1} Dependency-scoped repair for ch-a60b50ea300c848a: this fixed-k step does not rely on pending node 1.5.2.2.1. It depends directly on validated node 1.2, which proves beta lies in C. The root hypothesis for the chosen negative index k says that every x in C satisfies ||x - r\_k||\_1 <= D\_k. Since beta is in C by node 1.2, universal instantiation of that hypothesis at x = beta gives ||beta - r\_k||\_1 <= D\_k for this fixed k. \steptag{validated} \\[4pt]
\step{1.5.2.2.3} Because the preceding argument fixed only an arbitrary index k and used no property of k other than being one of the listed negative-part indices, the conclusion ||beta - r\_k||\_1 <= D\_k holds for every k. \steptag{validated} \\[6pt]
\step{1.5.2.2.3.1} Dependency-scoped repair for ch-96a8b58a8d27cb95: the universal conclusion is proved directly and does not use pending node 1.5.2.2.2. By validated node 1.2, beta lies in C. Now fix an arbitrary listed negative-part index k. The root hypothesis for D\_k says that every x in C satisfies ||x - r\_k||\_1 <= D\_k. Since beta lies in C, universal instantiation at x = beta gives ||beta - r\_k||\_1 <= D\_k for this arbitrary k. Therefore, by universal generalization over the listed indices k, ||beta - r\_k||\_1 <= D\_k holds for every k. \steptag{validated} \\[4pt]
\step{1.5.3} Multiplying the inequalities ||beta - r\_k||\_1 <= D\_k by c\_k >= 0 and summing over k gives sum\_k c\_k*||beta - r\_k||\_1 <= sum\_k c\_k*D\_k. \steptag{validated} \\[4pt]
\step{1.6} Since beta lies in C and m > 0, m * dist\_1(q, C) <= ||m*(q - beta)||\_1; combining the identity with the triangle inequality and the two residual bounds gives the claimed inequality. \steptag{validated} \\[6pt]
\step{1.6.1} Because beta lies in C, dist\_1(q, C) <= ||q - beta||\_1; because m > 0, multiplying gives m*dist\_1(q, C) <= m*||q - beta||\_1 = ||m*(q - beta)||\_1. \steptag{validated} \\[4pt]
\step{1.6.2} Using the identity from node 1.3 and the l1 triangle inequality gives ||m*(q - beta)||\_1 <= ||sum\_j b\_j*(p\_j - beta\_j)||\_1 + ||sum\_k c\_k*(beta - r\_k)||\_1. \steptag{validated} \\[6pt]
\step{1.6.2.1} Under the root hypotheses and the earlier choices B = sum\_j b\_j, G = sum\_k c\_k, beta = (sum\_j b\_j*beta\_j)/B, and m = B - G, the identity used here follows directly: multiplying q = (sum\_j b\_j*p\_j - sum\_k c\_k*r\_k)/m by m gives m*q = sum\_j b\_j*p\_j - sum\_k c\_k*r\_k, while m*beta = (B-G)*beta = sum\_j b\_j*beta\_j - sum\_k c\_k*beta; subtracting these two displayed equations gives m*(q - beta) = sum\_j b\_j*(p\_j - beta\_j) + sum\_k c\_k*(beta - r\_k). \steptag{validated} \\[4pt]
\step{1.6.2.2} Let A = sum\_j b\_j*(p\_j - beta\_j) and N = sum\_k c\_k*(beta - r\_k). By node 1.6.2.1, m*(q - beta) = A + N; therefore the l1 triangle inequality gives ||m*(q - beta)||\_1 = ||A + N||\_1 <= ||A||\_1 + ||N||\_1 = ||sum\_j b\_j*(p\_j - beta\_j)||\_1 + ||sum\_k c\_k*(beta - r\_k)||\_1. \steptag{validated} \\[4pt]
\step{1.6.3} The final inequality follows by a local chain: beta in C and m>0 give m*dist\_1(q,C) <= ||m*(q-beta)||\_1; the identity m*(q-beta)=sum\_j b\_j*(p\_j-beta\_j)+sum\_k c\_k*(beta-r\_k) and the l1 triangle inequality bound this by the sum of the positive and negative residual norms; those residual norms are bounded respectively by sum\_j b\_j*dist\_1(p\_j,C) and sum\_k c\_k*D\_k. Therefore m*dist\_1(q,C) <= sum\_j b\_j*dist\_1(p\_j,C) + sum\_k c\_k*D\_k. \steptag{validated} \\[6pt]
\step{1.6.3.1} Locally, because beta lies in C, dist\_1(q,C) <= ||q - beta||\_1; since m > 0, multiplying by m and using absolute homogeneity gives m*dist\_1(q,C) <= m*||q - beta||\_1 = ||m*(q - beta)||\_1. \steptag{validated} \\[4pt]
\step{1.6.3.2} Locally, the algebra identity m*(q - beta) = sum\_j b\_j*(p\_j - beta\_j) + sum\_k c\_k*(beta - r\_k), together with the l1 triangle inequality, gives ||m*(q - beta)||\_1 <= ||sum\_j b\_j*(p\_j - beta\_j)||\_1 + ||sum\_k c\_k*(beta - r\_k)||\_1. \steptag{validated} \\[6pt]
\step{1.6.3.2.1} In the local setup, q = (sum\_j b\_j*p\_j - sum\_k c\_k*r\_k)/m with m > 0, so multiplying by m gives m*q = sum\_j b\_j*p\_j - sum\_k c\_k*r\_k; subtracting m*beta from both sides gives m*(q - beta) = sum\_j b\_j*p\_j - sum\_k c\_k*r\_k - m*beta. \steptag{validated} \\[4pt]
\step{1.6.3.2.2} In the same local setup, write B = sum\_j b\_j and G = sum\_k c\_k. Since beta = (sum\_j b\_j*beta\_j)/B and m = B - G, we have B*beta = sum\_j b\_j*beta\_j and G*beta = sum\_k c\_k*beta, hence m*beta = (B - G)*beta = B*beta - G*beta = sum\_j b\_j*beta\_j - sum\_k c\_k*beta. \steptag{validated} \\[4pt]
\step{1.6.3.2.3} Using 1.6.3.2.1 and substituting the expression for m*beta from 1.6.3.2.2 gives m*(q - beta) = sum\_j b\_j*p\_j - sum\_k c\_k*r\_k - (sum\_j b\_j*beta\_j - sum\_k c\_k*beta); distributing the minus sign and grouping the j- and k-sums yields m*(q - beta) = sum\_j b\_j*(p\_j - beta\_j) + sum\_k c\_k*(beta - r\_k). \steptag{validated} \\[4pt]
\step{1.6.3.2.4} Let A = sum\_j b\_j*(p\_j - beta\_j) and N = sum\_k c\_k*(beta - r\_k). By 1.6.3.2.3, m*(q - beta) = A + N; applying the l1 triangle inequality gives ||m*(q - beta)||\_1 = ||A + N||\_1 <= ||A||\_1 + ||N||\_1 = ||sum\_j b\_j*(p\_j - beta\_j)||\_1 + ||sum\_k c\_k*(beta - r\_k)||\_1. \steptag{validated} \\[4pt]
\step{1.6.3.3} Locally, the positive residual obeys ||sum\_j b\_j*(p\_j - beta\_j)||\_1 <= sum\_j ||b\_j*(p\_j - beta\_j)||\_1 = sum\_j b\_j*||p\_j - beta\_j||\_1 = sum\_j b\_j*dist\_1(p\_j,C), using the l1 triangle inequality, b\_j >= 0, and the defining choice of beta\_j as a nearest point in C to p\_j. \steptag{validated} \\[4pt]
\step{1.6.3.4} Locally, the negative residual obeys ||sum\_k c\_k*(beta - r\_k)||\_1 <= sum\_k c\_k*||beta - r\_k||\_1 <= sum\_k c\_k*D\_k, using the l1 triangle inequality, c\_k >= 0, beta in C, and the hypothesis ||x-r\_k||\_1 <= D\_k for every x in C. \steptag{validated} \\[4pt]
\step{1.6.3.5} Combining the four local estimates 1.6.3.1--1.6.3.4 gives m*dist\_1(q,C) <= ||m*(q - beta)||\_1 <= ||sum\_j b\_j*(p\_j - beta\_j)||\_1 + ||sum\_k c\_k*(beta - r\_k)||\_1 <= sum\_j b\_j*dist\_1(p\_j,C) + sum\_k c\_k*D\_k, which is exactly the root inequality; thus the proof of 1.6.3 uses only these local children rather than pending sibling nodes 1.4, 1.5, 1.6.1, or 1.6.2. \steptag{archived} \\[4pt]
\step{1.6.3.6} Combining the four local estimates gives m*dist\_1(q,C) <= ||m*(q - beta)||\_1 <= ||sum\_j b\_j*(p\_j - beta\_j)||\_1 + ||sum\_k c\_k*(beta - r\_k)||\_1 <= sum\_j b\_j*dist\_1(p\_j,C) + sum\_k c\_k*D\_k, which is exactly the root inequality; this dependency-linked child is the local replacement for the former appeal to sibling nodes 1.4, 1.5, 1.6.1, and 1.6.2. \steptag{validated} \\[4pt]
\end{proofbox}

\end{document}
